\chapter{Abdominal Distension (Distended Abdomen)}

\section*{Definition}
Abdominal distension is a visible or measurable increase in abdominal size, often accompanied by pressure, fullness, or bloating. It may reflect gas, stool, fluid, pregnancy, an enlarged organ, or a mechanical obstruction. Distension is a physical finding; bloating is the subjective sensation of abdominal fullness, and the two may occur separately.

\section*{What Happens in Your Body}
Distension occurs when the abdominal cavity or gastrointestinal tract contains more volume than usual, or when the abdominal wall and diaphragm respond abnormally to normal intestinal contents. Gas, retained stool, impaired stomach or intestinal emptying, ascites, inflammation, and obstruction are important mechanisms. Severe distension can raise intra-abdominal pressure and impair venous return, breathing, or intestinal blood flow.

\section*{Causes}
\begin{itemize}
  \item Common functional causes include swallowed air, constipation, dietary fermentation, food intolerance, and irritable bowel syndrome.
  \item Gastrointestinal causes include ileus, intestinal obstruction, gastroparesis, inflammatory bowel disease, and severe gastroenteritis.
  \item Fluid accumulation (ascites) may result from cirrhosis, heart failure, cancer, or pancreatic disease.
  \item Pregnancy, obesity, enlarged organs, abdominal masses, and large ovarian or uterine lesions can increase abdominal size.
  \item Medication effects include constipation or slowed motility from opioids, anticholinergic medicines, calcium-channel blockers, and some iron preparations.
\end{itemize}

\section*{When to See a Doctor}
Seek emergency care for rapidly worsening distension with severe or continuous pain, repeated vomiting, a rigid or tender abdomen, inability to pass stool or gas, fever with toxicity, fainting, bloody or black stool, or breathing difficulty. Prompt assessment is also needed for new persistent distension, unexplained weight loss, jaundice, vomiting, pregnancy with pain, or distension that progressively increases.

\section*{Related Symptoms}
\begin{itemize}
  \item Abdominal pain, cramping, early satiety, nausea, or vomiting
  \item Constipation, diarrhea, inability to pass gas, or changes in stool caliber
  \item Shortness of breath, ankle swelling, jaundice, or easy bruising when fluid accumulation is present
\end{itemize}
\textbf{Important:} Distension with vomiting and inability to pass gas or stool can indicate intestinal obstruction and should not be treated with laxatives without medical advice.

\section*{Self-Care / Home Management}
\begin{itemize}
  \item For mild, intermittent symptoms without warning signs, eat slowly, avoid carbonated drinks, and record foods that reproducibly trigger symptoms.
  \item Maintain hydration, physical activity, and gradual dietary fiber intake; increase fiber slowly if constipation is present.
  \item Do not force food or fluids when vomiting or obstruction is suspected.
  \item Avoid stopping prescribed medicines or using repeated laxatives without discussing the cause with a clinician.
\end{itemize}

\section*{How Your Doctor Will Evaluate You}
\begin{itemize}
  \item History covers onset, progression, pain, bowel and urinary function, diet, medications, menstrual or pregnancy status, liver disease, and prior surgery.
  \item Examination assesses tympany, bowel sounds, tenderness, guarding, hernias, fluid wave, organ enlargement, and signs of dehydration or shock.
  \item Tests may include blood count, electrolytes, liver tests, lipase, urinalysis, pregnancy testing, and stool testing when indicated.
  \item Ultrasound or CT is selected for suspected ascites, mass, obstruction, inflammation, or an acute abdomen; plain radiography is used selectively.
\end{itemize}

\section*{Treatment Options}
Treatment addresses the cause and may include constipation management, dietary modification, treatment of infection or inflammation, drainage of clinically significant ascites, decompression for obstruction, or surgery for ischemia, perforation, or a structural lesion. Persistent unexplained distension warrants follow-up rather than repeated symptomatic treatment alone.

\section*{References}
\begin{thebibliography}{99}
\bibitem{abdominal_distension:ref1} Lacy BE, et al. ACG Clinical Guideline: Management of Irritable Bowel Syndrome. \textit{Am J Gastroenterol}. 2021;116(1):17--44.
\bibitem{abdominal_distension:ref2} American College of Radiology. ACR Appropriateness Criteria: Acute Nonlocalized Abdominal Pain. Revised 2023.
\bibitem{abdominal_distension:ref3} American Association for the Study of Liver Diseases. Diagnosis, Evaluation, and Management of Ascites. Hepatology guidance, 2021.
\end{thebibliography}
